\documentclass{article}
\usepackage{amsmath}
\begin{document}
TAME NODAL STACKY CURVES MARTIN BISHOP AND WILLIAM C. NEWMAN Abstract. In this paper we analyze the properties of tame nodal stacky curves, in particular twisted curves and doubly-twisted curves. Our main results are a complete classification of the possible struc- tures of a tame stacky node, along with computations of the Picard and Brauer groups of nodal stacky curves. 1. Introduction The geometry of stacky curves is controlled in large part by their Picard and Brauer groups. The Brauer group of a stacky curve is be- coming well understood through the recent work of Achenjang [Ach24] and Bishop [Bis25], and the Picard group of a smooth stacky curve was computed by Lopez [Lop23]. This paper, which originated as part of an upcoming work by the authors on twisted stable curves, aims to understand how the presence of stacky structure at nodes affects the Picard group. We will also finish a Brauer group computation which was promised in [Bis25, Example 3.9]. This paper gives an explicit description of the Picard and Brauer groups for tame nodal stacky curves, covering the two possible local structures at a node: twisted nodes, where the geometric stabilizer is cyclic and preserves the branches; and doubly-twisted , where there is a cyclic subgroup acting on each branch as well as a Z / 2 swapping branches. 1.1. Main results. At a node where µ n acts as ( ζ n , ζ a n ) for ζ n a prim- itive root of unity, set d − = gcd( a − 1 , n ) , d + = gcd( a + 1 , n ) , and n + = n/d + . Theorem (Theorem 5.7) . Let C be a stacky curve. Assume that C has a twisted node p . Let ˜ C be the partial normalization of C at p . Let ˜ C be the coarsening (see Definition 2.4) of ˜ C at the preimages of p . Let ˆ p denote either preimage. Then Pic C = k × ⊕ Pic ˜ C 1 n + ˆ p ⊕ Z /d + , 1 arXiv:2509.20629v1  [math.AG]  25 Sep 2025 2 MARTIN BISHOP AND WILLIAM C. NEWMAN where by Pic ˜ C D 1 n + ˆ p E we mean Pic ˜ C with an n th + root of O (ˆ p ) adjoined. The most common form of a twisted node is the balanced twisted node, where a ≡− 1 mod n . In fact, balanced twisted node is what most authors take twisted node to mean in recent work. Theorem 5.7 has a strong simplification in the balanced case. Theorem (Corollary 5.8) . Assume that C is as in Theorem 5.7. Fur- ther assume that the node p is balanced, that is, a ≡− 1 mod n . Then Pic C = k × ⊕ Pic ˜ C ⊕ Z /n. Theorem (Theorem 4.9) . Let C be a stacky curve. Assume that C has a doubly-twisted node p . Let C be its coarsening at p . Then Pic C = µ 2 ⊕ Pic C 1 d − p . Theorem (Proposition 4.14 and 5.5) . Let C be a stacky curve over an algebraically closed field k . The contribution of a twisted node to the Brauer group of C is H 2 ( G, k × triv ) = 0 . The contribution of a doubly- twisted node is H 2 ( G, k × triv ) = ( 0 if d − d + = n or G is non-split Z / 2 if d − d + = 2 n and G is split. These formulas track precisely how the introduction of stackiness alters the behavior of the classical, non-canonical splitting Pic C = k × ⊕ Pic ˜ C. They also contrast with the case of a smooth stacky curve. In [Lop23], Lopez shows that a smooth stacky curve is always obtained as a root stack from its coarse moduli space. This provides a uniform formula for computing Picard groups: it is always the pushout described in [Lop23, Theorem 1.1]. Whereas our Theorems 4.9 and 5.7 show that the Picard groups of nodal stacky curves, while captured in concise formulas, vary depending on the stack structure at the node and cannot be handled as uniformly. In particular, our results agree with the established fact that stacky nodes are not root stacks. 1.2. Conventions. We work over a field k , not assumed to be alge- braically closed except in the context of Brauer groups. All stacks are assumed to be tame, separated, Deligne-Mumford, and of finite type over k . Note that the tameness hypothesis means that char k is al- ways coprime to n and is coprime to 2 when discussing doubly-twisted curves. TAME NODAL STACKY CURVES 3 We take a stacky curve to be a tame, proper, geometrically integral, one-dimensional Deligne-Mumford stack with trivial generic stabilizer of finite type over a field. We assume that all of our nodes are split , meaning that each node as well as their preimages in a normalization are k -points. Throughout the paper, the group G will refer to either µ n or an extension of Z / 2 by µ n . We will call G split if G ∼ = µ n ⋊Z / 2. We also fix a 2 ≡ 1 mod n , d − = gcd( a − 1 , n ), d + = gcd( a + 1 , n ), and n + = n/d + . Whenever we discuss G , we are always in the situation of geometric points , so that µ n ∼ = Z /n non-canonically. 1.3. Acknowledgments. The first author would like to thank Gio- vanni Inchiostro, Cherry Ng, Brian Nugent, and Alex Schefflin for their helpful conversations. In particular, the authors would like to thank Jarod Alper and Yuchen Liu, who helped diagnose a persistent error with an early ver- sion of Theorem 5.7. The authors met at the 2023 AGNES Summer School on Intersec- tion Theory on Moduli Spaces. The authors thank the organizers and speakers. 2. Background 2.1. Assorted background. We recall Tsen’s Theorem for curves, as well as a stacky version. Theorem 2.1 (Tsen’s Theorem) . Suppose that C is a curve over an algebraically closed field. Then H 2 ( C, G m ) = 0 . Theorem 2.2. Let C be a stacky curve with trivial generic stabilizer over an algebraically closed field K . Let p 1 , . . . , p n be the singular points of C red and G 1 , . . . , G n be the geometric stabilizers at these points. Then H 2 ( C , G m ) = n M i =1 H 2 ( G i , K × ) where the action of each G i on K × is trivial. Proof. See [Bis25, Proposition 3.7]. □ Next, we describe the group cohomology of cyclic groups. Proposition 2.3. Let G be a cyclic group of order n and M a G - module. Let x denote the generator of G . Let N = 1 + x + x 2 + · · · + x n − 1 ∈ Z G. 4 MARTIN BISHOP AND WILLIAM C. NEWMAN Then H k ( G, M ) = ( M G /NM if k ≥ 2 even ker ( · N ) / ( x − 1) M if k odd . In particular, if the action of G on M is trivial then H k ( G, M ) = Hom( G, M ) if k is odd and H k ( G, M ) = 0 if k is even and multiplica- tion by | G | is surjective on M . Proof. See any standard reference for group cohomology, such as [DF04, Chapter 17] or [Wei94, Chapter 6]. □ 2.2. Extending a result of Meier. The main result of this section is Corollary 2.7, which was originally stated in the n = 2 case in [Mei18] and then generalized to the n ≥ 2 case in [Ach24] and [Bis25]. In [Ach24], conditions are given for when it applies for n ≥ 1. We now upgrade this to all cases for n ≥ 1, using Lemma 2.6. Definition 2.4 ([ATWo20], 2.3.1) . A morphism f : X → Y of stacks is called a coarsening if the inertia I X/Y is finite over X and for any flat Z → Y from an algebraic space Z the resulting base change X × Y Z → Z is a coarse moduli space. Proposition 2.5 ([ACV03] Proposition A.0.1, [Mei18] Proposition 8) . Let X be a separated tame Deligne-Mumford stack of finite type over k with coarsening f : X → X . Let x : Spec K → X be a geometric point. Let G = ( I X /X ) x be the relative stabilizer. Then for all n ≥ 0 ( R n f ∗ G m ) x = H n ( G, R × ) for some strictly Henselian local ring R with residue field K . Moreover, the group G acts trivially on K . Proof. See [Bis25, Proposition 2.11]. □ Lemma 2.6. Let A be a G -module. Let p be a prime that is coprime to | G | . Then for all n ≥ 1 , H n ( G, A ) = H n ( G, A [1 /p ]) . Proof. Let n ≥ 1. The result is trivial if p = 0, since we take A [1 /p ] to be A in this case. So assume p > 0. Let f : A → A [1 /p ] be the natural map. Consider the sequence 0 → ker f → A → A [1 /p ] → coker f → 0 . This leads to two short exact sequences (1) 0 → ker f → A → im f → 0 TAME NODAL STACKY CURVES 5 and (2) 0 → im f → A [1 /p ] → coker f → 0 . Taking cohomology for Equation 1 yields H n ( G, ker f ) → H n ( G, A ) → H n ( G, im f ) → H n +1 ( G, ker f ) . Note that f is injective on elements which are coprime to p , and hence ker f is a p -group. Therefore H i ( G, ker f ) = 0 for all i ≥ 1, since | G | is coprime to p . Thus H n ( G, A ) = H n ( G, im f ) . Repeating this for Equation 2 gives H n − 1 ( G, coker f ) → H n ( G, im f ) → H n ( G, A [1 /p ]) → H n ( G, coker f ) . Now, the cokernel of f is also a p -group, and hence H n ( G, coker f ) = 0 . Since H n ( G, im f ) contains no p -torsion, we see that the image of the p -group H n − 1 ( G, coker f ) is trivial (note that we are not saying that H n − 1 vanishes, since it may not in the case n = 1). Therefore H n ( G, A ) = H n ( G, im f ) = H n ( G, A [1 /p ]) . □ Corollary 2.7 ([Mei18] Lemma 9, [Ach24] Lemma 4.1, [Bis25] Corol- lary 2.12) . Let X and G be as in Proposition 2.5. Then for all geometric points x : Spec K → X and for all n ≥ 1 , ( R n f ∗ G m ) x = H n ( G, K × ) . Proof of Corollary 2.7. This proof is mostly an exact reproduction of the proof of Corollary 2.12 in [Bis25], but now extended to work for all n ≥ 1. If p = char K > 0, let p = char K ,and otherwise set p = 1. Using Proposition 2.5, all that’s left to show is that H n ( G, R × ) = H n ( G, K × ). Consider the exact sequence 1 → ker → R × [1 /p ] → K × [1 /p ] → 1 . Since R and K are both strictly Henselian, we have that R × [1 /p ] and K × [1 /p ] are both divisible. Therefore ker is divisible. Now by [Sta25, Tag 06RR], the torsion of R × [1 /p ] maps isomorphically onto the torsion of K × [1 /p ], and hence ker is torsion-free. The structure theorem for divisible groups (see [Fuc70, Theorem 23.1]) then implies that ker is a Q -vector space and hence has vanishing higher cohomology. Therefore, taking cohomology, we see that H n ( G, R × ) = H n ( G, R × [1 /p ]) = H n ( G, K × [1 /p ]) = H n ( G, K × ) 6 MARTIN BISHOP AND WILLIAM C. NEWMAN for all n ≥ 1, where the first and last equalities use Lemma 2.6 if p > 1. □ 3. Nodal stacky curves Definition 3.1. We define a nodal stacky curve C to be a proper, geometrically integral Deligne-Mumford stack of one dimension over a field k such that C has trivial generic stabilizer and C is covered with ´etale neighborhoods U →C such that U is (at worst) a nodal curve, in the classical sense. Equivalently, the completion of the strictly henselian local ring at every closed geometric point is either ¯ k [[ x ]] or ¯ k [[ x, y ]] / ( xy ). Remark 3.2. Recall from Section 1.2 that we only consider split nodes. The next result follows from and extends [Eke95], where its content is distributed throughout Section 1. Proposition 3.3. Let C be as above, and p ∈C a node. Let R denote the strict henselization of the local ring of k [ x, y ] / ( xy ) at the origin. Then C sh = C × C Spec O sh C,p is of the form [Spec R/G ] for G equal to one of either: • µ n , where µ n acts on ( x, y ) as ( ζ n x, ζ a n y ) for a primitive n th root of unity ζ n and an a ∈ ( Z /n ) × ; or • an extension of Z / 2 by µ n , where Z / 2 interchanges the axes and µ n acts on ( x, y ) as above with the additional condition that a 2 ≡ 1 mod n . Proof. The statement that C sh is of the form [Spec R/G ] for some finite G follows from the assumption that C has a node at p . Considering the action of G on the branches of C gives a map G → Z / 2, whose kernel consists of the branch-preserving elements of G . This kernel is necessarily cyclic, since the only geometric stabilizers of a smooth tame curve are cyclic. Therefore G is either cyclic (in the case where all elements are branch-preserving) or an extension of Z / 2 by µ n where Z / 2 acts by exchanging branches (in the case where there are branch-swapping elements). Asking that C have nontrivial geometric stabilizer requires that ζ n acts as (without loss of generality) ( ζ n x, ζ a n y ) for some primitive n th root of unity ζ n and a coprime to n , else the “ y -branch” will be stabilized everywhere by ζ n/ gcd( a,n ) n . In the case where G isn’t cyclic, we have the following further limitation on a . Let σ ∈ G denote a lift of the generator of Z / 2. Then σζ n σ − 1 · ( x, y ) = ( ζ a n x, ζ n y ) TAME NODAL STACKY CURVES 7 is of the form ( ζ k n x, ζ ak n y ) for some k . Hence k = a and ak = a 2 ≡ 1 mod n , as desired. □ Definition 3.4. In the case where G is an extension of Z / 2 by µ n , we say that C has a doubly-twisted node . Note 3.5. When all G are of the form µ n , C is called a twisted curve . Such curves have been studied extensively in the literature (see [AV02, AOV11, ACV03] among many others). While we allow a to take on any value relatively prime to n , most recent literature has focused on the balanced case where a ≡− 1 mod n . In fact, a ≡− 1 is becoming the standard usage of twisted . For non-cyclic geometric stabilizers, the particular case of G = µ n ⋊Z / 2 and a = − 1 (hence G ∼ = D n ) was discussed in [Rom05, Section 6], where such a node was said to be of dihedral type . 3.1. Numerology of doubly-twisted curves. For a doubly-twisted curve, the geometric stabilizer group G is by definition an extension of Z / 2 by µ n . We will work out some additional numerical properties of G beyond those of Proposition 3.3. We have already established that the action determines a ∈ Z /n such that a 2 ≡ 1 mod n . We can extract an additional piece of data from this. Let σ ∈ G be a lift of the generator of Z / 2. Then σ 2 ∈ µ n ≤ G , and so, fixing a primitive root ζ n ∈ µ n , we have σ 2 = ζ m n for some m . Notice that m is exactly the class which determines the extension class [ G ]. We will make this precise below. Proposition 3.6. For the action of Z / 2 on µ n determined by conju- gation, we have H 2 ( Z / 2 , µ n ) = ker(1 − a ) (1 + a ) µ n . Proof. This follows from Proposition 2.3. □ Proposition 3.7. The integer m is well-defined up to 1 + a , is in the kernel of multiplication by 1 − a , determines an element of H 2 ( Z / 2 , µ n ) , and determines the extension class of [ G ] . Moreover, the extension is trivial if and only if d + | m . Proof. Since σ 2 is fixed by conjugation, and conjugation acts as multi- plication by a , we see that m is in the kernel of multiplication by 1 − a in µ n , and is well-defined up to (1 + a ) µ n . Hence it determines a class in H 2 ( Z / 2 , µ n ), and this class corresponds precisely to [ G ]. To see the last statement, simply observe that the trivial extensions correspond to m ∈ (1 + a ) Z /n = d + Z /n. □ 8 MARTIN BISHOP AND WILLIAM C. NEWMAN Proposition 3.8. The group G is determined completely by a and m . Proof. The element a determines the action of Z / 2 on µ n . For a fixed action of Z / 2 on µ n , the extensions of Z / 2 by µ n with the given action are parametrized by H 2 ( Z / 2 , µ n ). By Proposition 3.7 m determines the class of [ G ] in H 2 . □ Definition 3.9. For the remainder of this section, we will refer to the extension of Z / 2 by µ n determined by ( a, m ) as G n,a,m . Proposition 3.10. We have that d − d + = n or 2 n . Proof. Let g = gcd( d − , d + ). Then g divides ( a + 1) − ( a − 1) = 2, and so g = 1 or 2. Since lcm( d − , d + ) = d − d + g divides n (as both d − and d + do), we see that d − d + divides gn , which is either n or 2 n . Because n divides a 2 − 1, we have n | d − d + . Hence d − d + is either n or 2 n . □ Finally, we show that any of these groups G do actually appear as the geometric stabilizer group at a node. Proposition 3.11. Let n ≥ 1 be an integer, a ∈ Z /n such that a 2 ≡ 1 mod n , and m ∈ Z /n such that m (1 − a ) ≡ 0 mod n . Then there exists a nodal stacky curve C with geometric stabilizer group G n,a,m (Definition 3.9) at the node. Proof. Consider the variety X , which is two copies of P 1 glued together at a point which affine locally looks like Spec( k [ x, y ] / ( xy )). On this open affine, we have automorphisms ( a, b ) : Spec( k [ x, y ] / ( xy )) → Spec( k [ x, y ] / ( xy )) ( x, y ) 7→ ( ax, by ) for a, b ∈ k × and τ : Spec( k [ x, y ] / ( xy )) → Spec( k [ x, y ] / ( xy )) ( x, y ) 7→ ( y, x ) . These automorphisms extend to the entirety of X . We consider the subgroup H of automorphisms generated by ( ζ n , ζ a n ) and τ (1 , ζ m n ). Ev- ery element of this subgroup is either of the form ( ζ k n , ζ ak n ) or τ ( ζ k n , ζ m + ak n ) TAME NODAL STACKY CURVES 9 for k ∈ Z /n . This group is isomorphic to G n,a,m and fixes the node, so [ X/H ] is a nodal stacky curve with geometric stabilizer group G n,a,m at the unique node. □ 4. A doubly twisted node Convention 4.1. Throughout this section we fix a tame stacky curve C over a field k which has a doubly-twisted node p : BG →C , as well as its coarsening at p , π : C → C . Let η : ˜ C →C be the partial normalization at p . Lastly, recall that we have fixed d + = gcd( a +1 , n ), d − = gcd( a − 1 , n ), and n + = n/d + . 4.1. The Leray spectral sequence. Following [Bis25], we use the Leray spectral sequence for π : C → C H p ( C, R q π ∗ G m ) = ⇒ H p + q ( C , G m ) . Since π ∗ G m = G m for coarsenings, the sequence of low-degree terms for this yields 0 → Pic( C ) → Pic( C ) → H 0 ( C, R 1 π ∗ G m ) → H 2 ( C, G m ) → H 2 ( C , G m ) , using that H 1 ( X , G m ) = Pic( X ). By Corollary 2.7, R 1 π ∗ G m vanishes away from the node, where its stalk is ( R 1 π ∗ G m ) p = H 1 ( G, ¯ k × triv ) with the trivial action of G on k × . Since H 1 ( G, ¯ k × triv ) = Hom( G, ¯ k × triv ) = Hom( G ab , ¯ k × ) the below sequence is exact 0 → Pic( C ) → Pic( C ) → Hom( G ab , ¯ k × )) → H 2 ( C, G m ) → H 2 ( C , G m ) . Lemma 4.2. The following sequence is exact (3) 0 → Pic( C ) → Pic( C ) → Hom( G ab , ¯ k × ) → 0 . Proof. It suffices to show that H 2 ( C, G m ) → H 2 ( C , G m ) is injective. Away from the node, we have C \ { p } = U = C \ { π ( p ) } . Consider the diagram H 2 ( C, G m ) H 2 ( C , G m ) H 2 ( U, G m ) . Since p is a smooth point of C (see the proof of Theorem 4.9 for an ex- planation of this), the restriction map to U is injective. Commutativity of the diagram forces the pullback to C to be injective as well. □ 10 MARTIN BISHOP AND WILLIAM C. NEWMAN Proposition 4.3. We have that G ab = ( Z / 2 ⊕ Z /d − , if both m and d − are even Z / 2 d − , else. Note that in either case, | G ab | = 2 d − . Note 4.4. All that we will really need is the statement on the size of G ab , but we include the full computation for completeness. Proof of Proposition 4.3. We know that G fits into the following exact sequence 1 → µ n → G → Z / 2 → 1 . Denote a generator of µ n by t and a lift of the generator of Z / 2 to G by σ . Then G has the following relations: t n = 1 σtσ − 1 = t a σ 2 = t m , subject to the numerical conditions from the beginning of Section 4. Abelianizing G adds the relation σt = tσ . If we consider the free abelian group Z ⟨ T ⟩⊕ Z ⟨ Σ ⟩ ∼ = Z 2 , then the above says that G ab is the quotient of this group by the relations • nT = 0, from t n = 1 in G , • ( a − 1) T = 0, from t = σtσ − 1 = t a in G ab , and • 2Σ − mT = 0, from σ 2 = t m in G . That is, G ab is the quotient of Z 2 by the Z -span of   n 0 a − 1 0 − m 2   . We may row-reduce the first two rows so that the top left entry is gcd( a − 1 , n ), i.e. d − . With this, G ab is now the quotient of Z 2 by the Z -span of  d − 0 − m 2  . Set k = gcd( d − , m, 2). Putting this matrix into Smith Normal Form shows that the span is precisely Z /k ⊕ Z / (2 d − /k ) = ( Z / 2 ⊕ Z /d − , if 2 divides both d − and m Z / (2 d − ) , else . □ Hence Pic C is an extension of a group of size 2 d − by Pic C , and all that remains is to determine which extension. TAME NODAL STACKY CURVES 11 4.2. Determining the extension. We will now use a classical argu- ment to compute Pic C . Definition 4.5. Define the sheaf Q on C as the following quotient: 0 →O × C → η ∗ O × ˜ C → Q → 0 . Since pushforward along a finite morphism is exact in the ´etale topol- ogy, the above short exact sequence gives the long exact sequence 0 →O C ( C ) × →O ˜ C ( ˜ C ) × → H 0 ( C , Q ) → Pic C → Pic ˜ C . But C and ˜ C are both integral and proper, and hence O C ( C ) × = O ˜ C ( ˜ C ) × = k × . Therefore we have the following exact sequence: (4) 0 → H 0 ( C , Q ) → Pic C → Pic ˜ C . We now begin an analysis of the terms of Equation 4. Proposition 4.6. Let k × inv denote the sheaf on BG associated to the G -module whose underlying group is k × , where the action is given by inversion from Z / 2 . Then Q ∼ = p ∗ k × inv . Proof. By definition Q ( U ) is the sheafification of the assignment U 7→ ( η ∗ G m )( U ) G m ( U ) = G m ( ˜ U ) G m ( U ) . This is now reduced to the level of curves, where we know that this quotient is, Zariski locally, just ( k × ) m where m is the number of preim- ages of the node. But the assignment U 7→ ( k × ) m is a familiar sheaf: namely, it is the pushforward of k × from the residual gerbe at the node, as desired. Now we must determine the action of G . We note that the k × appearing in ( k × ) m is actually k × × k × k × where the quotient is by the diagonal ( t, t ). Since the branch-preserving part of G acts trivially on k , it acts trivially on the quotient. However, the branch-swapping part of G does not act trivially. More specifically, the Z / 2 interchanges the two copies of k × and hence acts by inversion on the quotient, as claimed. □ Corollary 4.7. The global sections of Q are H 0 ( C , Q ) = µ 2 . Proof. Exactness of pushforward under a closed immersion together with Proposition 4.6 implies that H 0 ( C , Q ) = H 0 ( BG, k × inv ) = H 0 ( G, k × inv ) = µ 2 . 12 MARTIN BISHOP AND WILLIAM C. NEWMAN Where the last equality holds since inversion fixes µ 2 . Note that tame- ness exclude characteristic two, so µ 2 is nontrivial. □ We have then shown that the following sequence is exact (5) 0 → µ 2 → Pic C → Pic ˜ C . Notice that, using Equation 3, the inclusion of µ 2 into Pic C is split by the map Pic C → Hom( G ab , k × ). This proves the following: Proposition 4.8. We have Pic C = µ 2 ⊕ im h Pic C → Pic ˜ C i . With this, we can finish the computation of the Picard group. Theorem 4.9. Let C be a stacky curve with coarse moduli space C . Assume that C has a doubly-twisted node p . Let C be its coarsening at p . Then Pic C = µ 2 ⊕ Pic C 1 d − p , where by Pic C D 1 d − p E we mean Pic C with a d th − root of O ( p ) adjoined. Proof. By Proposition 4.8, we only need to compute im h Pic C → Pic ˜ C i . Since the ´etale local structure of C at the node is  Spec k [ x, y ] / ( xy ) G  where G swaps the branches, the ´etale local structure of the normal- ization of C is  Spec k [ z ] µ n  , i.e. it is smooth with geometric stabilizer µ n . Moreover, the composi- tion ˜ C →C → C is also the coarse moduli space of ˜ C . From [Lop23] we then have an isomorphism Pic ˜ C ∼ = Pic C 1 n p . By Equation 3 and Proposition 4.3, we have that Pic( C ) , → Pic( C ) and the cokernel has size 2 d − . Then, we must have that the cokernel of Pic( C ) , → Pic( C ) /µ 2 has size d − . As Pic( C ) ⟨ 1 n p ⟩ / Pic( C ) ∼ = Z /n , we see that the image of Pic( C ) must be Pic C D 1 d − p E . □ TAME NODAL STACKY CURVES 13 Example 4.10. From Theorem 4.9, any choice of an action with a = 1 leads to the largest possible Picard group for a nodal curve. When a = 1, there are at most two such extensions: Z / 2 ⊕ Z /n and Z / 2 n . We consider the interesting case of Z / 2 n . Fix a primitive 2 n th root of unity ζ 2 n . Let G = Z / 2 n act as 1 · ( x, y ) = ( ζ 2 n y, ζ 2 n x ). Then this is a doubly-twisted node whose geometric sta- bilizer is cyclic, and Pic C = µ 2 ⊕ Pic C 1 n p . 4.3. The Brauer group of C . In this section, all group actions on k × are taken to be trivial. We also assume that k is algebraically closed. Recall that G is an extension of Z / 2 by µ n : 1 → µ n → G → Z / 2 → 1 . We will use the Lyndon-Hochschild-Serre spectral sequence E p,q 2 = H p ( Z / 2 , H q ( µ n , k × )) = ⇒ H p + q ( G, k × ) to compute H 2 ( G, k × ) . Lemma 4.11. We have H 2 ( G, k × ) = ker  d 1 , 1 2 : H 1 ( Z / 2 , H 1 ( µ n , k × )) → H 3 ( Z / 2 , k × )  . Proof. The E 2 , 0 2 and E 0 , 2 2 terms are E 2 , 0 2 = H 2 ( Z / 2 , k × ) = k × / ( k × ) 2 = 0 and E 0 , 2 2 = H 0 ( Z / 2 , H 2 ( µ n , k × )) = H 0 ( Z / 2 , 0) = 0 . Therefore the end of the sequence of low-degree terms reads 0 → H 2 ( G, k × ) → H 1 ( Z / 2 , H 1 ( µ n , k × )) d 1 , 1 2 −−→ H 3 ( Z / 2 , k × ) , and so H 2 ( G, k × ) = ker( d 1 , 1 2 : E 1 , 1 2 → E 3 , 0 2 ) as desired. □ Lemma 4.12. The E 1 , 1 2 term is H 1 ( Z / 2 , H 1 ( µ n , k × )) = Z /d + Z / ( n/d − ) , which is either trivial if d − d + = n or Z / 2 if d − d + = 2 n . The generator of this group is n/d + . 14 MARTIN BISHOP AND WILLIAM C. NEWMAN Proof. The formulas in Proposition 2.3 tell us that H 1 ( Z / 2 , H 1 ( µ n , k × )) = H 1 ( Z / 2 , Z /n ) where the action of Z / 2 on Z /n is multiplication by a . Further appli- cation of these formulas gives H 1 ( Z / 2 , Z /n ) = { k ∈ Z /n : ( a + 1) k = 0 } ( a − 1) Z /n = Z /d + Z / ( n/d − ) . This group is either trivial if d − d + = n or Z / 2 if d − d + = 2 n , and its generator in either case is the element n/d + . □ Note 4.13. In Lemma 4.12 we identified H 1 ( µ n , k × ) = Hom( µ n , k × ) = Z /n. Unravelling this identification to write in terms of Hom( µ n , k × ), which will be necessary in the proof of Theorem 4.14, we see that the generator is  φ : ζ n 7→ ζ n/d + n  . Proposition 4.14. Suppose that C has a doubly-twisted node over an algebraically closed field. The contribution of the node to the Brauer group of C is H 2 ( G, k × triv ) = ( 0 if d − d + = n or G is non-split Z / 2 if d − d + = 2 n and G is split. Proof. By Lemmas 4.11 and 4.12, we may assume that d − d + = 2 n , else E 1 , 1 2 and hence H 2 ( G, k × ) is trivial. By Note 4.13, the generator of H 1 ( Z / 2 , H 1 ( µ n , k × )) is the element  φ : ζ n 7→ ζ n/d + n  . Recall (Section 3.1) that the class of the extension [ G ] is determined by m . By [HS53, Theorem 4], for the element φ ∈ H 1 ( Z / 2 , H 1 ( µ n , k × )), we have d 1 , 1 2 ( φ ) = φ ( ζ − m n ) ∈ µ 2 = Hom( Z / 2 , k × ) = H 3 ( Z / 2 , k × ) . This shows that if the extension is trivial (i.e. split), then d 1 , 1 2 is trivial and hence H 2 ( G, k × ) = H 1 ( Z / 2 , Z /n ) = Z / 2 . In the non-split case, we consider d 1 , 1 2 ( φ ) = φ ( ζ − m n ) = ζ − mn/d + n . TAME NODAL STACKY CURVES 15 Note that this is equal to 1 if and only if d + | m , which is precisely the criterion for whether or not the extension is trivial from Proposition 3.7. Hence d 1 , 1 2 ( φ ) ̸ = 1 when d − d + = 2 n and the extension is nontrivial. We conclude that d 1 , 1 2 is injective and H 2 ( G, k × ) = 0. □ Example 4.15. Theorem 4.14 now provides the example promised in [Bis25, Example 3.9]. Pick any split extension G = µ n ⋊Z / 2 with a chosen so that d − d + = 2 n . For k algebraically closed, set C =  Spec k [ x, y ] / ( xy ) G  where G acts as ( ζ n , ζ a n ) from µ n and swaps branches from Z / 2. Then H 2 ( C , G m ) = H 2 ( G, k × ) = Z / 2 . For an example of such an a , we can pick a = − 1. Hence G = D n , and so C is a balanced doubly-twisted curve (or is of dihedral type). Any such a and G provide, from the perspective of complexity of the singularity, the simplest possible example of a stacky curve over an algebraically closed field with non-trivial Brauer group. 5. A twisted node Convention 5.1. In this section we fix a stacky nodal curve C with a twisted node, p . Denote the geometric stabilizer at p , which is µ n , by G . Let π : C → C be the coarsening at p , let η : ˜ C →C be the partial normalization of C at p , and let ˜ C be the coarsening of ˜ C at the preimages of the node (which is also the partial normalization of C ). Lastly, recall that we have fixed d + = gcd( a +1 , n ), d − = gcd( a − 1 , n ), and n + = n/d + . Lemma 5.2. The below sequence is exact: 0 → Pic C → Pic C → Z /n → 0 . Proof. The beginning of the sequence of low-degree terms of the Leray spectral sequence for π : C → C reads 0 → Pic C → Pic C → H 0 ( C, R 1 π ∗ G m ) → H 2 ( C, G m ) → H 2 ( C , G m ) . As in Section 4, R 1 π ∗ G m vanishes except at the node, where its stalk is Z /n . Therefore it suffices to show the injectivity of H 2 ( C, G m ) → H 2 ( C , G m ) . 16 MARTIN BISHOP AND WILLIAM C. NEWMAN Following Lemma 4.2, away from the node we have C \ { π ( p ) } = U = C \ { p } . Consider the diagram H 2 ( C, G m ) H 2 ( C , G m ) H 2 ( U, G m ) . To show that H 2 ( C, G m ) → H 2 ( C , G m ) is injective, it suffices to show that the restriction of H 2 ( C, G m ) to U is injective. However, since the node is split, the kernel of this restriction is precisely ker " Br k → 2 M i =1 Br k # , where the map into each direct summand is the identity. Therefore the restriction and hence pullback is injective. □ Following Section 4, let Q be the quotient 0 →O × C → η ∗ O × ˜ C → Q → 0 . As before, we have the exact sequence 0 → H 0 ( C , Q ) → Pic C → Pic ˜ C → H 1 ( C , Q ) → H 2 ( C , G m ) → H 2 ( ˜ C , G m ) . Proposition 5.3. Let k × triv denote the sheaf on BG associated to the G -module whose underlying group is k × , where the action is trivial. Then Q ∼ = p ∗ k × triv . Proof. This is identical to Proposition 4.6, except now the action of G on k × is trivial. □ Corollary 5.4. The cohomology of Q is H 0 ( C , Q ) = k × and H 1 ( C , Q ) = H 1 ( G, k × triv ) . Proof. This is also identical to its Section 4 counterpart, now with trivial action. Hence H 0 ( C , Q ) = H 0 ( G, k × ) = k × and H 1 ( C , Q ) = H 1 ( G, k × ) = Hom( µ n , k × ) = Z /n. □ Proposition 5.5. Let C be a stacky curve over an algebraically closed field with a twisted node. The contribution of the node to H 2 ( C , G m ) is H 2 ( G, k × triv ) = 0 . TAME NODAL STACKY CURVES 17 Proof. This follows from Proposition 2.3, since cyclic groups have van- ishing second cohomology with coefficients in k × . □ Proposition 5.6. The pullback H 2 ( C , G m ) → H 2 ( ˜ C , G m ) is injective. Proof. In fact a stronger statement is true. Since H 2 ( µ n , k × triv ) = 0, we see that twisted nodes do not contribute to the Brauer group of a stacky curve. Therefore since C and ˜ C are isomorphic away from the twisted node and its preimages (which also have cyclic stabilizers), and since the preimages are both k -points, the pullback is an isomorphism. □ Hence the following sequence is exact (6) 0 → k × → Pic C → Pic ˜ C → Z /n → 0 . Theorem 5.7. Let C be a stacky curve. Assume that C has a twisted node p . Let ˜ C be the partial normalization of C at p . Let ˜ C be the coarsening of ˜ C at the preimages of p . Let ˆ p denote either preimage. Then Pic C = k × ⊕ Pic ˜ C 1 n + ˆ p ⊕ Z /d + . Proof. From Equation 6 we have Pic C = k × ⊕ im(Pic C → Pic ˜ C ) , so we need to determine the image of the pullback. Let p ∈C be the node and p 1 , p 2 its preimages in ˜ C . By [Lop23], the Picard group of ˜ C is Pic ˜ C = Pic ˜ C 1 n p 1 , 1 n p 2 ∼ = Pic ˜ C 1 n ˆ p ⊕ Z /n, where ˆ p can be chosen to be either preimage of p in ˜ C and the Z /n direct summand consists of line bundles corresponding to k ( 1 n p 1 − 1 n p 2 ). Therefore any direct summand appearing in Pic C will be precisely the intersection of the image of Pic C with the subgroup of line bundles of the form O ( k ( 1 n p 1 − 1 n p 2 )). Now, using the Leray spectral sequence to write 0 → Pic C → Pic C → Z /n → 0 , pick a set-theoretic lift L of 1 ∈ Z /n such that its component in Pic C is trivial. It follows from the ´etale local description of C that the pullback of L to ˜ C is 1 n p 1 + a 1 n p 2 . The intersection of the span of this line bundle with those of the form O ( k ( 1 n p 1 − 1 n p 2 )) is isomorphic to Z /d + . Hence Pic C = k × ⊕ Pic ˜ C 1 n + ˆ p ⊕ Z /d + . □ 18 MARTIN BISHOP AND WILLIAM C. NEWMAN Corollary 5.8. Assume that C is as in Theorem 5.7. Further assume that the node p is balanced, that is, a ≡− 1 mod n . Then Pic C = k × ⊕ Pic ˜ C ⊕ Z /n. TAME NODAL STACKY CURVES 19 References [Ach24] Niven Achenjang. On brauer groups of tame stacks, 2024. [ACV03] Dan Abramovich, Alessio Corti, and Angelo Vistoli. Twisted bundles and admissible covers. volume 31, pages 3547–3618. 2003. Special issue in honor of Steven L. Kleiman. [AOV11] Dan Abramovich, Martin Olsson, and Angelo Vistoli. Twisted stable maps to tame Artin stacks. J. Algebraic Geom. , 20(3):399–477, 2011. [ATWo20] Dan Abramovich, Michael Temkin, and Jaros l aw W l odarczyk. Toroidal orbifolds, destackification, and Kummer blowings up. Algebra Number Theory , 14(8):2001–2035, 2020. With an appendix by David Rydh. [AV02] Dan Abramovich and Angelo Vistoli. Compactifying the space of stable maps. J. Amer. Math. Soc. , 15(1):27–75, 2002. [Bis25] Martin Bishop. Brauer groups of tame stacky curves and their µ r -gerbes, 2025. [DF04] David S. Dummit and Richard M. Foote. Abstract algebra . John Wiley & Sons, Inc., Hoboken, NJ, third edition, 2004. [DHW12] Karel Dekimpe, Manfred Hartl, and Sarah Wauters. A seven-term exact sequence for the cohomology of a group extension. J. Algebra , 369:70–95, 2012. [Eke95] Torsten Ekedahl. Boundary behaviour of Hurwitz schemes. In The mod- uli space of curves (Texel Island, 1994) , volume 129 of Progr. Math. , pages 173–198. Birkh¨auser Boston, Boston, MA, 1995. [Fuc70] L´aszl´o Fuchs. Infinite abelian groups. Vol. I , volume Vol. 36 of Pure and Applied Mathematics . Academic Press, New York-London, 1970. [HS53] G. Hochschild and J.-P. Serre. Cohomology of group extensions. Trans. Amer. Math. Soc. , 74:110–134, 1953. [Lop23] Rose Lopez. Picard groups of stacky curves. 2023. [Mei18] Lennart Meier. Computing brauer groups using coarse moduli – draft version. 2018. [Rom05] Matthieu Romagny. Group actions on stacks and applications. Michigan Math. J. , 53(1):209–236, 2005. [Sta25] The Stacks Project Authors. Stacks Project . https://stacks.math. columbia.edu , 2025. [Wei94] Charles A. Weibel. An introduction to homological algebra , volume 38 of Cambridge Studies in Advanced Mathematics . Cambridge University Press, Cambridge, 1994. 
\end{document}